\Chapter{102}

\Verse{1}
{
    \zR{话说王夫人打发人来唤宝钗，宝钗连忙过来请了安。}
}
{
    \eR{[English source not present in the checked-in Bencraft/Joly files.]}
}
{
}

\Verse{2}
{
    \zR{王夫人道：“你三妹妹如 今要出嫁了，你们作嫂子的大家开导开导他，也是你们姊妹之情。}
}
{
    \eR{[English source not present in the checked-in Bencraft/Joly files.]}
}
{
}

\Verse{3}
{
    \zR{况且他也是个明 白孩子，我看你们两个也很合的来。}
}
{
    \eR{[English source not present in the checked-in Bencraft/Joly files.]}
}
{
}

\Verse{4}
{
    \zR{只是我听见说，宝玉听见他三妹妹出门子，哭 的了不的。}
}
{
    \eR{[English source not present in the checked-in Bencraft/Joly files.]}
}
{
}

\Verse{5}
{
    \zR{你也该劝劝他才是。}
}
{
    \eR{[English source not present in the checked-in Bencraft/Joly files.]}
}
{
}

\Verse{6}
{
    \zR{如今我的身子是十病九痛的，你二嫂子也是三日好 两日不好。}
}
{
    \eR{[English source not present in the checked-in Bencraft/Joly files.]}
}
{
}

\Verse{7}
{
    \zR{你还心地明白些，诸事该管的，也别说只管吞着，不肯得罪人。}
}
{
    \eR{[English source not present in the checked-in Bencraft/Joly files.]}
}
{
}

\Verse{8}
{
    \zR{将来这 一番家事都是你的担子。”}
}
{
    \eR{[English source not present in the checked-in Bencraft/Joly files.]}
}
{
}

\Verse{9}
{
    \zR{宝钗答应着。}
}
{
    \eR{[English source not present in the checked-in Bencraft/Joly files.]}
}
{
}

\Verse{10}
{
    \zR{王夫人又说道：“还有一件事，你二嫂子昨 儿带了柳家媳妇的丫头来，说补在你们屋里。”}
}
{
    \eR{[English source not present in the checked-in Bencraft/Joly files.]}
}
{
}

\Verse{11}
{
    \zR{宝钗道：“今日平儿才带过来，说是 太太和二奶奶的主意。”}
}
{
    \eR{[English source not present in the checked-in Bencraft/Joly files.]}
}
{
}

\Verse{12}
{
    \zR{王夫人道：“是呦，你二嫂子和我说，我想也没要紧，不便 驳他的回。}
}
{
    \eR{[English source not present in the checked-in Bencraft/Joly files.]}
}
{
}

\Verse{13}
{
    \zR{只是一件，我见那孩子眉眼儿上头也不是个很安顿的。}
}
{
    \eR{[English source not present in the checked-in Bencraft/Joly files.]}
}
{
}

\Verse{14}
{
    \zR{起先为宝玉房里 的丫头狐狸似的，我撵了几个，那时候你也自然知道，才搬回家去的。}
}
{
    \eR{[English source not present in the checked-in Bencraft/Joly files.]}
}
{
}

\Verse{15}
{
    \zR{如今有你， 固然不比先前了。}
}
{
    \eR{[English source not present in the checked-in Bencraft/Joly files.]}
}
{
}

\Verse{16}
{
    \zR{我告诉你，不过留点神儿就是了。}
}
{
    \eR{[English source not present in the checked-in Bencraft/Joly files.]}
}
{
}

\Verse{17}
{
    \zR{你们屋里，就是袭人那孩子还 可以使得。”}
}
{
    \eR{[English source not present in the checked-in Bencraft/Joly files.]}
}
{
}

\Verse{18}
{
    \zR{宝钗答应了，又说了几句话，便过来了。}
}
{
    \eR{[English source not present in the checked-in Bencraft/Joly files.]}
}
{
}

\Verse{19}
{
    \zR{饭后到了探春那边，自有一 番殷勤劝慰之言，不必细说。}
}
{
    \eR{[English source not present in the checked-in Bencraft/Joly files.]}
}
{
}

\Verse{20}
{
    \zR{次日，探春将要起身，又来辞宝玉。}
}
{
    \eR{[English source not present in the checked-in Bencraft/Joly files.]}
}
{
}

\Verse{21}
{
    \zR{宝玉自然难割难分。}
}
{
    \eR{[English source not present in the checked-in Bencraft/Joly files.]}
}
{
}

\Verse{22}
{
    \zR{探春倒将纲常大体的 话，说的宝玉始而低头不语，后来转悲作喜，似有醒悟之意。}
}
{
    \eR{[English source not present in the checked-in Bencraft/Joly files.]}
}
{
}

\Verse{23}
{
    \zR{于是探春放心辞别众 人，竟上轿登程，水舟陆车而去。}
}
{
    \eR{[English source not present in the checked-in Bencraft/Joly files.]}
}
{
}

\Verse{24}
{
    \zR{先前众姊妹们都住在大观园中，后来贾妃薨后，也不修葺。}
}
{
    \eR{[English source not present in the checked-in Bencraft/Joly files.]}
}
{
}

\Verse{25}
{
    \zR{到了宝玉娶亲，林 黛玉一死，史湘云回去，宝琴在家住着，园中人少，况兼天气寒冷，李纨姊妹、探 春、惜春等俱挪回旧所。}
}
{
    \eR{[English source not present in the checked-in Bencraft/Joly files.]}
}
{
}

\Verse{26}
{
    \zR{到了花朝月夕，依旧相约玩耍。}
}
{
    \eR{[English source not present in the checked-in Bencraft/Joly files.]}
}
{
}

\Verse{27}
{
    \zR{如今探春一去，宝玉病后 不出屋门，益发没有高兴的人了。}
}
{
    \eR{[English source not present in the checked-in Bencraft/Joly files.]}
}
{
}

\Verse{28}
{
    \zR{所以园中寂寞，只有几家看园的人住着。}
}
{
    \eR{[English source not present in the checked-in Bencraft/Joly files.]}
}
{
}

\Verse{29}
{
    \zR{那日，尤氏过来送探春起身，因天晚省得套车，便从前年在园里开通宁府的那 个便门里走过去了。}
}
{
    \eR{[English source not present in the checked-in Bencraft/Joly files.]}
}
{
}

\Verse{30}
{
    \zR{觉得凄凉满目，台榭依然，女墙一带都种作园地一般，心中怅 然如有所失。}
}
{
    \eR{[English source not present in the checked-in Bencraft/Joly files.]}
}
{
}

\Verse{31}
{
    \zR{因到家中，便有些身上发热。}
}
{
    \eR{[English source not present in the checked-in Bencraft/Joly files.]}
}
{
}

\Verse{32}
{
    \zR{扎挣一两天，竟躺倒了。}
}
{
    \eR{[English source not present in the checked-in Bencraft/Joly files.]}
}
{
}

\Verse{33}
{
    \zR{日间的发烧犹 可，夜里身热异常，便谵语绵绵。}
}
{
    \eR{[English source not present in the checked-in Bencraft/Joly files.]}
}
{
}

\Verse{34}
{
    \zR{贾珍连忙请了大夫看视，说感冒起的，如今缠经 入了足阳明胃经，所以谵语不清，如有所见，有了大秽即可身安。}
}
{
    \eR{[English source not present in the checked-in Bencraft/Joly files.]}
}
{
}

\Verse{35}
{
    \zR{尤氏服了两剂， 并不稍减，更加发起狂来。}
}
{
    \eR{[English source not present in the checked-in Bencraft/Joly files.]}
}
{
}

\Verse{36}
{
    \zR{贾珍着急，便叫贾蓉来：“打听外头有好医生，再请几 位来瞧瞧。”}
}
{
    \eR{[English source not present in the checked-in Bencraft/Joly files.]}
}
{
}

\Verse{37}
{
    \zR{贾蓉回道：“前儿这个大夫是最兴时的了，只怕我母亲的病不是药治得 好的。”}
}
{
    \eR{[English source not present in the checked-in Bencraft/Joly files.]}
}
{
}

\Verse{38}
{
    \zR{贾珍道：“胡说，不吃药，难道由他去罢？”}
}
{
    \eR{[English source not present in the checked-in Bencraft/Joly files.]}
}
{
}

\Verse{39}
{
    \zR{贾蓉道：“不是说不治，为的 是前日母亲往西府去，回来是穿着园子里走过来的。}
}
{
    \eR{[English source not present in the checked-in Bencraft/Joly files.]}
}
{
}

\Verse{40}
{
    \zR{一到了家就身上发烧，别是撞 客着了罢。}
}
{
    \eR{[English source not present in the checked-in Bencraft/Joly files.]}
}
{
}

\Verse{41}
{
    \zR{外头有个毛半仙，是南方人，卦起的很灵，不如请他来占算占算。}
}
{
    \eR{[English source not present in the checked-in Bencraft/Joly files.]}
}
{
}

\Verse{42}
{
    \zR{看有 信儿呢，就依着他；要是不中用，再请别的好大夫来。”}
}
{
    \eR{[English source not present in the checked-in Bencraft/Joly files.]}
}
{
}

\Verse{43}
{
    \zR{贾珍听了，即刻叫人请来；坐在书房内喝了茶，便说：“府上叫我，不知占什 么事？”}
}
{
    \eR{[English source not present in the checked-in Bencraft/Joly files.]}
}
{
}

\Verse{44}
{
    \zR{贾蓉道：“家母有病，请教一卦。”}
}
{
    \eR{[English source not present in the checked-in Bencraft/Joly files.]}
}
{
}

\Verse{45}
{
    \zR{毛半仙道：“既如此，取净水洗手，设 下香案，让我起出一课来看就是了。”}
}
{
    \eR{[English source not present in the checked-in Bencraft/Joly files.]}
}
{
}

\Verse{46}
{
    \zR{一时，下人安排定了，他便怀里掏出卦筒来， 走到上头，恭恭敬敬的作了一个揖，手内摇着卦筒，口里念道：“伏以太极两仪， ？？}
}
{
    \eR{[English source not present in the checked-in Bencraft/Joly files.]}
}
{
}

\Verse{47}
{
    \zR{交感，图书出而变化不穷，神圣作而诚求必应。}
}
{
    \eR{[English source not present in the checked-in Bencraft/Joly files.]}
}
{
}

\Verse{48}
{
    \zR{兹有信官贾某，为因母病，虔 请伏羲、文王、周公、孔子四大圣人，鉴临在上，诚感则灵，有凶报凶，有吉报吉。}
}
{
    \eR{[English source not present in the checked-in Bencraft/Joly files.]}
}
{
}

\Verse{49}
{
    \zR{先请内象三爻。”}
}
{
    \eR{[English source not present in the checked-in Bencraft/Joly files.]}
}
{
}

\Verse{50}
{
    \zR{说着，将筒内的钱倒在盘内，说：“有灵的，头一爻就是‘交’。”}
}
{
    \eR{[English source not present in the checked-in Bencraft/Joly files.]}
}
{
}

\Verse{51}
{
    \zR{拿起来又摇了一摇，倒出来，说是“单”。}
}
{
    \eR{[English source not present in the checked-in Bencraft/Joly files.]}
}
{
}

\Verse{52}
{
    \zR{第三爻又是“交”。}
}
{
    \eR{[English source not present in the checked-in Bencraft/Joly files.]}
}
{
}

\Verse{53}
{
    \zR{检起钱来，嘴里说是： “内爻已示，更请外象三爻，完成一卦。”}
}
{
    \eR{[English source not present in the checked-in Bencraft/Joly files.]}
}
{
}

\Verse{54}
{
    \zR{起出来，是“单拆单”。}
}
{
    \eR{[English source not present in the checked-in Bencraft/Joly files.]}
}
{
}

\Verse{55}
{
    \zR{那毛半仙收了卦 筒和铜钱，便坐下问道：“请坐，请坐，让我来细细的看看。}
}
{
    \eR{[English source not present in the checked-in Bencraft/Joly files.]}
}
{
}

\Verse{56}
{
    \zR{这个卦乃是‘未济’ 之卦。}
}
{
    \eR{[English source not present in the checked-in Bencraft/Joly files.]}
}
{
}

\Verse{57}
{
    \zR{世爻是第三爻，午火兄弟劫财，晦气是一定该有的。}
}
{
    \eR{[English source not present in the checked-in Bencraft/Joly files.]}
}
{
}

\Verse{58}
{
    \zR{如今尊驾为母问病，用 神是初爻，真是父母爻动出官鬼来。}
}
{
    \eR{[English source not present in the checked-in Bencraft/Joly files.]}
}
{
}

\Verse{59}
{
    \zR{五爻上又有一层官鬼，我看令堂太夫人的病是 不轻的。}
}
{
    \eR{[English source not present in the checked-in Bencraft/Joly files.]}
}
{
}

\Verse{60}
{
    \zR{还好，还好，如今子亥之水休囚，寅木动而生火。}
}
{
    \eR{[English source not present in the checked-in Bencraft/Joly files.]}
}
{
}

\Verse{61}
{
    \zR{世爻上动出一个子孙来， 倒是克鬼的。}
}
{
    \eR{[English source not present in the checked-in Bencraft/Joly files.]}
}
{
}

\Verse{62}
{
    \zR{况且日月生身，再隔两日，子水官鬼落空，交到戌日就好了。}
}
{
    \eR{[English source not present in the checked-in Bencraft/Joly files.]}
}
{
}
\ChapterImage{img/chapters/205第102-103回 寧國府骨肉病災祲 大觀園符水驅妖孽.jpg}


\Verse{63}
{
    \zR{但是父 母爻上变鬼，恐怕令尊大人也有些关碍。}
}
{
    \eR{[English source not present in the checked-in Bencraft/Joly files.]}
}
{
}

\Verse{64}
{
    \zR{就是本身世爻比劫过重，到了水旺土衰的 日子也不好。”}
}
{
    \eR{[English source not present in the checked-in Bencraft/Joly files.]}
}
{
}

\Verse{65}
{
    \zR{说完了，便撅着胡子坐着。}
}
{
    \eR{[English source not present in the checked-in Bencraft/Joly files.]}
}
{
}

\Verse{66}
{
    \zR{贾蓉起先听他捣鬼，心里忍不住要笑；听他讲的卦理明白，又说生怕父亲也不 好，便说道：“卦是极高明的，但不知我母亲到底是什么病？”}
}
{
    \eR{[English source not present in the checked-in Bencraft/Joly files.]}
}
{
}

\Verse{67}
{
    \zR{毛半仙道：“据这卦 上，世爻午火变水相克，必是寒火凝结。}
}
{
    \eR{[English source not present in the checked-in Bencraft/Joly files.]}
}
{
}

\Verse{68}
{
    \zR{若要断得清楚，揲蓍也不大明白，除非用 ‘大六壬’才断的准。”}
}
{
    \eR{[English source not present in the checked-in Bencraft/Joly files.]}
}
{
}

\Verse{69}
{
    \zR{贾蓉道：“先生都高明的么？”}
}
{
    \eR{[English source not present in the checked-in Bencraft/Joly files.]}
}
{
}

\Verse{70}
{
    \zR{毛半仙道：“知道些。”}
}
{
    \eR{[English source not present in the checked-in Bencraft/Joly files.]}
}
{
}

\Verse{71}
{
    \zR{贾蓉 便要请教，报了一个时辰。}
}
{
    \eR{[English source not present in the checked-in Bencraft/Joly files.]}
}
{
}

\Verse{72}
{
    \zR{毛先生便画了盘子，将神将排定算去，是戌上白虎。}
}
{
    \eR{[English source not present in the checked-in Bencraft/Joly files.]}
}
{
}

\Verse{73}
{
    \zR{“这 课加做‘魄化课’。}
}
{
    \eR{[English source not present in the checked-in Bencraft/Joly files.]}
}
{
}

\Verse{74}
{
    \zR{大凡白虎乃是凶将，乘旺象气受制，便不能为害。}
}
{
    \eR{[English source not present in the checked-in Bencraft/Joly files.]}
}
{
}

\Verse{75}
{
    \zR{如今乘着死 神死煞及时令囚死，则为饿虎，定是伤人。}
}
{
    \eR{[English source not present in the checked-in Bencraft/Joly files.]}
}
{
}

\Verse{76}
{
    \zR{就如魄神受惊消散，故名‘魄化’。}
}
{
    \eR{[English source not present in the checked-in Bencraft/Joly files.]}
}
{
}

\Verse{77}
{
    \zR{这 课象说是人身丧魄，忧患相仍，病多丧死，讼有忧惊。}
}
{
    \eR{[English source not present in the checked-in Bencraft/Joly files.]}
}
{
}

\Verse{78}
{
    \zR{按象有日暮虎临，必定是傍 晚得病的。}
}
{
    \eR{[English source not present in the checked-in Bencraft/Joly files.]}
}
{
}

\Verse{79}
{
    \zR{象内说：‘凡占此课，必定旧宅有伏虎作怪，或有形响。’}
}
{
    \eR{[English source not present in the checked-in Bencraft/Joly files.]}
}
{
}

\Verse{80}
{
    \zR{如今尊驾为大 人而占，正合着虎在阳忧男，在阴忧女，此课十分凶险呢。”}
}
{
    \eR{[English source not present in the checked-in Bencraft/Joly files.]}
}
{
}

\Verse{81}
{
    \zR{贾蓉没有听完，唬得 面上失色道：“先生说的很是，但与那卦又不大相合，到底有妨碍么？”}
}
{
    \eR{[English source not present in the checked-in Bencraft/Joly files.]}
}
{
}

\Verse{82}
{
    \zR{毛半仙道： “你不用慌，待我慢慢的再看。”}
}
{
    \eR{[English source not present in the checked-in Bencraft/Joly files.]}
}
{
}

\Verse{83}
{
    \zR{低着头又咕哝了一会子，便说：“好了，有救星了。}
}
{
    \eR{[English source not present in the checked-in Bencraft/Joly files.]}
}
{
}

\Verse{84}
{
    \zR{算出巳上有贵神救解，谓之‘魄化魂归’，先忧后喜，是不妨事的，只要小心些就 是了。”}
}
{
    \eR{[English source not present in the checked-in Bencraft/Joly files.]}
}
{
}

\Verse{85}
{
    \zR{贾蓉奉上卦金，送了出去，回禀贾珍，说是：“母亲的病，是在旧宅傍晚得的， 为撞着什么‘伏尸白虎’。”}
}
{
    \eR{[English source not present in the checked-in Bencraft/Joly files.]}
}
{
}

\Verse{86}
{
    \zR{贾珍道：“你说你母亲前日从园里走回来的，可不是那 里撞着的!你还记得你二婶娘到园里去，回来就病了？}
}
{
    \eR{[English source not present in the checked-in Bencraft/Joly files.]}
}
{
}

\Verse{87}
{
    \zR{他虽没有见什么，后来那些丫 头老婆们都说是山子上一个毛烘烘的东西，眼睛有灯笼大，还会说话，他把二奶奶 赶回来了，唬出一场病来。”}
}
{
    \eR{[English source not present in the checked-in Bencraft/Joly files.]}
}
{
}

\Verse{88}
{
    \zR{贾蓉道：“怎么不记得!我还听见宝二叔家的焙茗说： 晴雯做了园里芙蓉花的神了；林姑娘死了，半空里有音乐，必定他也是管什么花儿 了。}
}
{
    \eR{[English source not present in the checked-in Bencraft/Joly files.]}
}
{
}

\Verse{89}
{
    \zR{想这许多妖怪在园里，还了得。}
}
{
    \eR{[English source not present in the checked-in Bencraft/Joly files.]}
}
{
}

\Verse{90}
{
    \zR{头里人多阳气重，常来常往不打紧；如今冷落 的时候，母亲打那里走，还不知踹了什么花儿呢，不然就是撞着那一个。}
}
{
    \eR{[English source not present in the checked-in Bencraft/Joly files.]}
}
{
}

\Verse{91}
{
    \zR{那卦也还 算是准的。”}
}
{
    \eR{[English source not present in the checked-in Bencraft/Joly files.]}
}
{
}

\Verse{92}
{
    \zR{贾珍道：“到底说有妨碍没有呢？”}
}
{
    \eR{[English source not present in the checked-in Bencraft/Joly files.]}
}
{
}

\Verse{93}
{
    \zR{贾蓉道：“据他说，到了戌日就好 了。}
}
{
    \eR{[English source not present in the checked-in Bencraft/Joly files.]}
}
{
}

\Verse{94}
{
    \zR{只愿早两天好，或除两天才好。”}
}
{
    \eR{[English source not present in the checked-in Bencraft/Joly files.]}
}
{
}

\Verse{95}
{
    \zR{贾珍道：“这又是什么意思？”}
}
{
    \eR{[English source not present in the checked-in Bencraft/Joly files.]}
}
{
}

\Verse{96}
{
    \zR{贾蓉道：“那 先生若是这样准，生怕老爷也有些不自在。”}
}
{
    \eR{[English source not present in the checked-in Bencraft/Joly files.]}
}
{
}

\Verse{97}
{
    \zR{正说着，里头喊说：“奶奶要坐起到那 边园里去，丫头们都按捺不住。”}
}
{
    \eR{[English source not present in the checked-in Bencraft/Joly files.]}
}
{
}

\Verse{98}
{
    \zR{贾珍等进去安慰，只闻尤氏嘴里乱说：“穿红的来 叫我!穿绿的来赶我！”}
}
{
    \eR{[English source not present in the checked-in Bencraft/Joly files.]}
}
{
}

\Verse{99}
{
    \zR{地下这些人又怕又好笑。}
}
{
    \eR{[English source not present in the checked-in Bencraft/Joly files.]}
}
{
}

\Verse{100}
{
    \zR{贾珍便命人买些纸钱，送到园里烧 化。}
}
{
    \eR{[English source not present in the checked-in Bencraft/Joly files.]}
}
{
}

\Verse{101}
{
    \zR{果然那夜出了汗，便安静些。}
}
{
    \eR{[English source not present in the checked-in Bencraft/Joly files.]}
}
{
}

\Verse{102}
{
    \zR{到了戌日，也就渐渐的好起来。}
}
{
    \eR{[English source not present in the checked-in Bencraft/Joly files.]}
}
{
}

\Verse{103}
{
    \zR{由是，一人传十，十人传百，都说大观园中有了妖怪，唬得那些看园的人也不 修花补树、灌溉果蔬。}
}
{
    \eR{[English source not present in the checked-in Bencraft/Joly files.]}
}
{
}

\Verse{104}
{
    \zR{起先晚上不敢行走，以致鸟兽逼人；近来甚至日间也是约伴 持械而行。}
}
{
    \eR{[English source not present in the checked-in Bencraft/Joly files.]}
}
{
}

\Verse{105}
{
    \zR{过了些时，果然贾珍也病，竟不请医调治，轻则到园化纸许愿，重则详 星拜斗。}
}
{
    \eR{[English source not present in the checked-in Bencraft/Joly files.]}
}
{
}

\Verse{106}
{
    \zR{贾珍方好，贾蓉等相继而病。}
}
{
    \eR{[English source not present in the checked-in Bencraft/Joly files.]}
}
{
}

\Verse{107}
{
    \zR{如此接连数月，闹的两府俱怕。}
}
{
    \eR{[English source not present in the checked-in Bencraft/Joly files.]}
}
{
}

\Verse{108}
{
    \zR{从此风声鹤 唳，草木皆妖。}
}
{
    \eR{[English source not present in the checked-in Bencraft/Joly files.]}
}
{
}

\Verse{109}
{
    \zR{园中出息一概全蠲，各房月例重新添起，反弄的荣府中更加拮据。}
}
{
    \eR{[English source not present in the checked-in Bencraft/Joly files.]}
}
{
}

\Verse{110}
{
    \zR{那些看园的没有了想头，个个要离此处，每每造言生事，便将花妖树怪编派起来， 各要搬出，将园门封固，再无人敢到园中。}
}
{
    \eR{[English source not present in the checked-in Bencraft/Joly files.]}
}
{
}

\Verse{111}
{
    \zR{以致崇楼高阁，琼馆瑶台，皆为禽兽所 栖。}
}
{
    \eR{[English source not present in the checked-in Bencraft/Joly files.]}
}
{
}

\Verse{112}
{
    \zR{却说晴雯的表兄吴贵正住在园门口。}
}
{
    \eR{[English source not present in the checked-in Bencraft/Joly files.]}
}
{
}

\Verse{113}
{
    \zR{他媳妇自从晴雯死后，听见说作了花神， 每日晚间便不敢出门。}
}
{
    \eR{[English source not present in the checked-in Bencraft/Joly files.]}
}
{
}

\Verse{114}
{
    \zR{这一日吴贵出门买东西，回来晚了。}
}
{
    \eR{[English source not present in the checked-in Bencraft/Joly files.]}
}
{
}

\Verse{115}
{
    \zR{那媳妇子本有些感冒着 了，日间吃错了药，晚上吴贵到家，已死在炕上。}
}
{
    \eR{[English source not present in the checked-in Bencraft/Joly files.]}
}
{
}

\Verse{116}
{
    \zR{外面的人因那媳妇子不大妥当， 便说妖怪爬过墙来吸了精去死的。}
}
{
    \eR{[English source not present in the checked-in Bencraft/Joly files.]}
}
{
}

\Verse{117}
{
    \zR{于是老太太着急的了不得，另派了好些人将宝玉 的住房围住，巡逻打更。}
}
{
    \eR{[English source not present in the checked-in Bencraft/Joly files.]}
}
{
}

\Verse{118}
{
    \zR{这些小丫头们还说，有看见红脸的，有看见很俊的女人的， 吵嚷不休，唬的宝玉天天害怕。}
}
{
    \eR{[English source not present in the checked-in Bencraft/Joly files.]}
}
{
}

\Verse{119}
{
    \zR{亏得宝钗有把持，听见丫头们混说，便吓唬着要打， 所以那些谣言略好些。}
}
{
    \eR{[English source not present in the checked-in Bencraft/Joly files.]}
}
{
}

\Verse{120}
{
    \zR{无奈各房的人都是疑人疑鬼的不安静，也添了人坐更，于是 更加了好些食用。}
}
{
    \eR{[English source not present in the checked-in Bencraft/Joly files.]}
}
{
}

\Verse{121}
{
    \zR{独有贾赦不大很信，说：“好好儿的园子，那里有什么鬼怪。”}
}
{
    \eR{[English source not present in the checked-in Bencraft/Joly files.]}
}
{
}

\Verse{122}
{
    \zR{挑了个风清日暖 的日子，带了好几个家人，手内持着器械，到园踹看动静。}
}
{
    \eR{[English source not present in the checked-in Bencraft/Joly files.]}
}
{
}

\Verse{123}
{
    \zR{众人劝他不依。}
}
{
    \eR{[English source not present in the checked-in Bencraft/Joly files.]}
}
{
}

\Verse{124}
{
    \zR{到了园 中，果然阴气逼人。}
}
{
    \eR{[English source not present in the checked-in Bencraft/Joly files.]}
}
{
}
\ChapterImage{img/chapters/206第102-103回 施-毒-計-金-桂-自-焚-身(原簽誤)..jpg}


\Verse{125}
{
    \zR{贾赦还扎挣前走，跟的人都探头缩脑的。}
}
{
    \eR{[English source not present in the checked-in Bencraft/Joly files.]}
}
{
}

\Verse{126}
{
    \zR{内中有个年轻的家人， 心内已经害怕，只听唿的一声，回过头来，只见五色灿烂的一件东西跳过去了，唬 的“嗳哟”一声，腿子发软，就躺倒了。}
}
{
    \eR{[English source not present in the checked-in Bencraft/Joly files.]}
}
{
}

\Verse{127}
{
    \zR{贾赦回身查问，那小子喘嘘嘘的回道：“亲 眼看见一个黄脸红胡子绿衣裳一个妖精!走到树林子后头山窟窿里去了。”}
}
{
    \eR{[English source not present in the checked-in Bencraft/Joly files.]}
}
{
}

\Verse{128}
{
    \zR{贾赦听 了，便也有些胆怯，问道：“你们都看见么？”}
}
{
    \eR{[English source not present in the checked-in Bencraft/Joly files.]}
}
{
}

\Verse{129}
{
    \zR{有几个推顺水船儿的回说：“怎么没 瞧见？}
}
{
    \eR{[English source not present in the checked-in Bencraft/Joly files.]}
}
{
}

\Verse{130}
{
    \zR{因老爷在头里，不敢惊动罢了。}
}
{
    \eR{[English source not present in the checked-in Bencraft/Joly files.]}
}
{
}

\Verse{131}
{
    \zR{奴才们还掌得住。”}
}
{
    \eR{[English source not present in the checked-in Bencraft/Joly files.]}
}
{
}

\Verse{132}
{
    \zR{说得贾赦害怕，也不敢再 走。}
}
{
    \eR{[English source not present in the checked-in Bencraft/Joly files.]}
}
{
}

\Verse{133}
{
    \zR{急急的回来，吩咐小子们：“不用提及，只说看遍了，没有什么东西。”}
}
{
    \eR{[English source not present in the checked-in Bencraft/Joly files.]}
}
{
}

\Verse{134}
{
    \zR{心里实 也相信，要到真人府里请法官驱邪。}
}
{
    \eR{[English source not present in the checked-in Bencraft/Joly files.]}
}
{
}

\Verse{135}
{
    \zR{岂知那些家人无事还要生事，今见贾赦怕了， 不但不瞒着，反添些穿凿，说得人人吐舌。}
}
{
    \eR{[English source not present in the checked-in Bencraft/Joly files.]}
}
{
}

\Verse{136}
{
    \zR{贾赦没法，只得请道士到园作法，驱邪 逐妖。}
}
{
    \eR{[English source not present in the checked-in Bencraft/Joly files.]}
}
{
}

\Verse{137}
{
    \zR{择吉日，先在省亲正殿上铺排起坛场来。}
}
{
    \eR{[English source not present in the checked-in Bencraft/Joly files.]}
}
{
}

\Verse{138}
{
    \zR{供上三清圣像，旁设二十八宿并马、 赵、温、周四大将，下排三十六天将图像。}
}
{
    \eR{[English source not present in the checked-in Bencraft/Joly files.]}
}
{
}

\Verse{139}
{
    \zR{香花灯烛设满一堂，钟鼓法器排列两边， 插着五方旗号。}
}
{
    \eR{[English source not present in the checked-in Bencraft/Joly files.]}
}
{
}

\Verse{140}
{
    \zR{道纪司派定四十九位道众的执事，净了一天坛。}
}
{
    \eR{[English source not present in the checked-in Bencraft/Joly files.]}
}
{
}

\Verse{141}
{
    \zR{三位法官行香取水 毕，然后擂起法鼓。}
}
{
    \eR{[English source not present in the checked-in Bencraft/Joly files.]}
}
{
}

\Verse{142}
{
    \zR{法师们俱戴上七星冠，披上九宫八卦的法衣，踏着登云履，手 执牙笏，便拜表请圣。}
}
{
    \eR{[English source not present in the checked-in Bencraft/Joly files.]}
}
{
}

\Verse{143}
{
    \zR{又念了一天的消灾驱邪接福的《洞玄经》，以后便出榜召将。}
}
{
    \eR{[English source not present in the checked-in Bencraft/Joly files.]}
}
{
}

\Verse{144}
{
    \zR{榜上大书“太乙、混元、上清三境灵宝符？}
}
{
    \eR{[English source not present in the checked-in Bencraft/Joly files.]}
}
{
}

\Verse{145}
{
    \zR{演教大法师，行文敕令本境诸神到坛听 用”。}
}
{
    \eR{[English source not present in the checked-in Bencraft/Joly files.]}
}
{
}

\Verse{146}
{
    \zR{那日两府上下爷们仗着法师擒妖，都到园中观看，都说：“好大法令，呼神遣 将的闹起来，不管有多少妖怪也唬跑了。”}
}
{
    \eR{[English source not present in the checked-in Bencraft/Joly files.]}
}
{
}

\Verse{147}
{
    \zR{大家都挤到坛前。}
}
{
    \eR{[English source not present in the checked-in Bencraft/Joly files.]}
}
{
}

\Verse{148}
{
    \zR{只见小道士们将旗幡 举起，按定五方站住，伺候法师号令。}
}
{
    \eR{[English source not present in the checked-in Bencraft/Joly files.]}
}
{
}

\Verse{149}
{
    \zR{三位法师，一位手提宝剑，拿着法水，一位 捧着七星皂旗，一位举着桃木打妖鞭，立在坛前。}
}
{
    \eR{[English source not present in the checked-in Bencraft/Joly files.]}
}
{
}

\Verse{150}
{
    \zR{只听法器一停，上头令牌三下， 口中念起咒来，那五方旗便团团散布。}
}
{
    \eR{[English source not present in the checked-in Bencraft/Joly files.]}
}
{
}

\Verse{151}
{
    \zR{法师下坛，叫本家领着到各处楼阁殿亭，房 廊屋舍，山崖水畔，洒了法水，将剑指画了一回。}
}
{
    \eR{[English source not present in the checked-in Bencraft/Joly files.]}
}
{
}

\Verse{152}
{
    \zR{回来，连击令牌，将七星旗祭起， 众道士将旗幡一聚接下，打妖鞭望空打了三下。}
}
{
    \eR{[English source not present in the checked-in Bencraft/Joly files.]}
}
{
}

\Verse{153}
{
    \zR{本家众人都道拿住妖怪，争着要看， 及到跟前，并不见有什么形响。}
}
{
    \eR{[English source not present in the checked-in Bencraft/Joly files.]}
}
{
}

\Verse{154}
{
    \zR{只见法师叫众道士拿取瓶罐，将妖收下，加上封条， 法师朱笔书符收起，令人带回在本观塔下镇住，一面撤坛谢将。}
}
{
    \eR{[English source not present in the checked-in Bencraft/Joly files.]}
}
{
}

\Verse{155}
{
    \zR{贾赦恭敬叩谢了法 师。}
}
{
    \eR{[English source not present in the checked-in Bencraft/Joly files.]}
}
{
}

\Verse{156}
{
    \zR{贾蓉等小弟兄背地都笑个不住，说：“这样的大排场，我打量拿着妖怪，给我 们瞧瞧到底是些什么东西，那里知道是这样搜罗。}
}
{
    \eR{[English source not present in the checked-in Bencraft/Joly files.]}
}
{
}

\Verse{157}
{
    \zR{究竟妖怪拿去了没有？”}
}
{
    \eR{[English source not present in the checked-in Bencraft/Joly files.]}
}
{
}

\Verse{158}
{
    \zR{贾珍听 见，骂道：“糊涂东西!妖怪原是聚则成形，散则成气，如今多少神将在这里，还敢 现形吗？}
}
{
    \eR{[English source not present in the checked-in Bencraft/Joly files.]}
}
{
}

\Verse{159}
{
    \zR{无非把这妖气收了，便不作祟，就是法力了。”}
}
{
    \eR{[English source not present in the checked-in Bencraft/Joly files.]}
}
{
}

\Verse{160}
{
    \zR{众人将信将疑，且等不见响 动再说。}
}
{
    \eR{[English source not present in the checked-in Bencraft/Joly files.]}
}
{
}

\Verse{161}
{
    \zR{那些下人只知妖怪被擒，疑心去了，便不大惊小怪，往后果然没人提起了。}
}
{
    \eR{[English source not present in the checked-in Bencraft/Joly files.]}
}
{
}

\Verse{162}
{
    \zR{贾 珍等病愈复原，都道法师神力。}
}
{
    \eR{[English source not present in the checked-in Bencraft/Joly files.]}
}
{
}

\Verse{163}
{
    \zR{独有一个小厮笑说道：“头里那些响动，我也不知 道。}
}
{
    \eR{[English source not present in the checked-in Bencraft/Joly files.]}
}
{
}

\Verse{164}
{
    \zR{就是跟着大老爷进园这一日，明明是个大公野鸡飞过去了。}
}
{
    \eR{[English source not present in the checked-in Bencraft/Joly files.]}
}
{
}

\Verse{165}
{
    \zR{拴儿吓离了眼，说 的活像，我们都替他圆了个谎，大老爷就认真起来。}
}
{
    \eR{[English source not present in the checked-in Bencraft/Joly files.]}
}
{
}

\Verse{166}
{
    \zR{倒瞧了个很热闹的坛场。”}
}
{
    \eR{[English source not present in the checked-in Bencraft/Joly files.]}
}
{
}

\Verse{167}
{
    \zR{众 人虽然听见，那里肯信，究无人敢住。}
}
{
    \eR{[English source not present in the checked-in Bencraft/Joly files.]}
}
{
}

\Verse{168}
{
    \zR{一日，贾赦无事，正想要叫几个家下人搬住园中看守，惟恐夜晚藏匿奸人。}
}
{
    \eR{[English source not present in the checked-in Bencraft/Joly files.]}
}
{
}

\Verse{169}
{
    \zR{方 欲传出话去，只见贾琏进来，请了安，回说：“今日到大舅家去，听见一个荒信， 说是二叔被节度使参进来，为的是失察属员，重征粮米，请旨革职的事。”}
}
{
    \eR{[English source not present in the checked-in Bencraft/Joly files.]}
}
{
}

\Verse{170}
{
    \zR{贾赦听 了，吃惊道：“只怕是谣言罢？}
}
{
    \eR{[English source not present in the checked-in Bencraft/Joly files.]}
}
{
}

\Verse{171}
{
    \zR{前儿你二叔带书子来说，探春于某日到了任所，择了 某日吉时，送了你妹子到了海疆，路上风恬浪静，合家不必挂念。}
}
{
    \eR{[English source not present in the checked-in Bencraft/Joly files.]}
}
{
}

\Verse{172}
{
    \zR{还说节度认亲， 倒设席贺喜。}
}
{
    \eR{[English source not present in the checked-in Bencraft/Joly files.]}
}
{
}

\Verse{173}
{
    \zR{那里有做了亲戚倒提参起来的？}
}
{
    \eR{[English source not present in the checked-in Bencraft/Joly files.]}
}
{
}

\Verse{174}
{
    \zR{且不必言语，快到吏部打听明白，就 来回我。”}
}
{
    \eR{[English source not present in the checked-in Bencraft/Joly files.]}
}
{
}

\Verse{175}
{
    \zR{贾琏即刻出去，不到半日回来，便说：“才到吏部打听，果然二叔被参。}
}
{
    \eR{[English source not present in the checked-in Bencraft/Joly files.]}
}
{
}

\Verse{176}
{
    \zR{题本上去，亏得皇上的恩典，没有交部，便下旨意，说是：‘失察属员，重征粮米， 苛虐百姓，本应革职，姑念初膺外任，不谙吏治，被属员蒙蔽，着降三级，加恩仍
        以工部员外上行走，并令即日回京。’}
}
{
    \eR{[English source not present in the checked-in Bencraft/Joly files.]}
}
{
}

\Verse{177}
{
    \zR{这信是准的。}
}
{
    \eR{[English source not present in the checked-in Bencraft/Joly files.]}
}
{
}

\Verse{178}
{
    \zR{正在吏部说话的时候，来了一 个江西引见的知县，说起我们二叔是很感激的。}
}
{
    \eR{[English source not present in the checked-in Bencraft/Joly files.]}
}
{
}

\Verse{179}
{
    \zR{但说是个好上司，只是用人不当， 那些家人在外招摇撞骗，欺凌属员，已经把好名声都弄坏了。}
}
{
    \eR{[English source not present in the checked-in Bencraft/Joly files.]}
}
{
}

\Verse{180}
{
    \zR{节度大人早已知道， 也说我们二叔是个好人。}
}
{
    \eR{[English source not present in the checked-in Bencraft/Joly files.]}
}
{
}

\Verse{181}
{
    \zR{不知怎么样，这回又参了。}
}
{
    \eR{[English source not present in the checked-in Bencraft/Joly files.]}
}
{
}

\Verse{182}
{
    \zR{想是忒闹得不好，恐将来弄出 大祸，所以借了一件失察的事情参的，倒是避重就轻的意思，也未可知。”}
}
{
    \eR{[English source not present in the checked-in Bencraft/Joly files.]}
}
{
}

\Verse{183}
{
    \zR{贾赦未 听说完，便叫贾琏：“先去告诉你婶子知道，且不必告诉老太太就是了。”}
}
{
    \eR{[English source not present in the checked-in Bencraft/Joly files.]}
}
{
}

\Verse{184}
{
    \zR{贾琏去回 王夫人。}
}
{
    \eR{[English source not present in the checked-in Bencraft/Joly files.]}
}
{
}

\Verse{185}
{
    \zR{未知有何话说，下回分解。}
}
{
    \eR{[English source not present in the checked-in Bencraft/Joly files.]}
}
{
}

\Verse{186}
{
    \zR{(本章完)}
}
{
    \eR{[English source not present in the checked-in Bencraft/Joly files.]}
}
{
}
